\chapter*{Préface}
\addcontentsline{toc}{chapter}{Préface}
\markboth{Préface}{Préface}

\section*{À qui ce manuel s'adresse}

Ce manuel s'adresse à la personne qui installe, configure et exploite
une instance mybibli pour son foyer, une petite association ou une
bibliothèque communautaire. Il suppose que vous savez :

\begin{itemize}
  \item ouvrir un terminal et lancer des commandes
        \texttt{docker compose} ;
  \item modifier un fichier texte de configuration dans l'éditeur de
        votre choix ;
  \item lire une URL imprimée dans une ligne de log et l'ouvrir dans
        un navigateur.
\end{itemize}

Aucune connaissance de Rust, MariaDB ou HTMX n'est requise. Le manuel
se limite délibérément au rôle d'\emph{exploitant} — compiler
mybibli depuis les sources ou modifier son code est couvert par
\texttt{CLAUDE.md} et les artefacts de planification du dépôt.

\section*{Conventions}

\textbf{Commandes.} Les fragments de shell sont présentés dans un
encadré à largeur fixe :

\begin{lstlisting}
docker compose up -d
\end{lstlisting}

\textbf{Fichiers de configuration.} Les extraits de fichiers utilisent
le même style, le nom du fichier étant indiqué dans le texte
environnant. Les commentaires à l'intérieur du fragment utilisent le
marqueur de commentaire propre au fichier.

\textbf{Encadrés.} Trois types d'encadrés signalent les passages
importants :

\begin{note}
Une \emph{note} fournit un contexte ou un éclairage qui aide à
comprendre pourquoi quelque chose fonctionne comme cela. La sauter
ne pose pas de problème.
\end{note}

\begin{tip}
Une \emph{astuce} suggère une pratique qui rend l'usage quotidien
plus fluide. Essayez-la une fois que les bases sont maîtrisées.
\end{tip}

\begin{warning}
Un \emph{avertissement} signale une étape où une erreur peut
entraîner une perte de données, vous verrouiller hors du panneau
d'administration ou exposer l'instance à des accès non autorisés.
Lisez ces encadrés en entier.
\end{warning}

\textbf{Renvois croisés.} Les chapitres et sections sont liés dans le
PDF — cliquez sur un titre ou un numéro de page en bleu pour y
sauter. L'\hyperref[index]{index} en fin d'ouvrage renvoie chaque
terme important à sa page de premier usage.

\section*{Demander de l'aide}

La source de vérité canonique pour mybibli est le dépôt GitHub :

\begin{center}
\url{https://github.com/guycorbaz/mybibli}
\end{center}

Ouvrez une \emph{issue} sur le dépôt \texttt{guycorbaz/mybibli}
lorsque vous rencontrez un bug, une lacune dans la documentation ou
une fonctionnalité manquante. Les modèles d'issue demandent la
version (visible en pied de page du panneau d'administration), votre
méthode d'installation et les étapes pour reproduire le problème.

\section*{Couverture du manuel}

Ce manuel couvre mybibli \textbf{1.1.0 et versions ultérieures}. Les
versions antérieures comportaient un bug de sécurité où l'assistant
de première mise en service était silencieusement contourné — voir le
chapitre~\ref{ch:upgrade} pour le chemin de mise à jour depuis 1.0.x.

\vspace{1cm}
\hrule
\vspace{0.5cm}
\textit{Bon catalogage.}
